\chapter{\texttt{tricycle.py}: the car model}
\label{ch:tricycle_py}

\textbf{Role:} defines the car as an MJCF string and offers a helper to write
it to disk for conversion. At run time only three constants are imported from
it (\code{CHASSIS_Z}, \code{JOINT_STEER}, \code{JOINT_REAR}); the MJCF itself
is used offline (Chapter~\ref{ch:asset}). 156 lines.

\section{Header}

\codepart{\srcroot/luna_hifi_tasks/tricycle/tricycle.py}{1}{29}

\begin{explain}
\lis{1}{4} What the file is.
\lis{6}{9} The most important operational fact about this file: it is
  \emph{not} read during training. The USD under \srcpath{assets/} is, and it
  must be regenerated when this file changes.
\lis{11}{18} An ASCII plan view of the body so that the numbers below have a
  picture: the two rails meet at the apex over the front wheel; the crossbar
  closes the rear; the deck sits inside.
\lis{20}{23} Why the wheels are spheres. The implicit MPM solver treats
  colliders as signed-distance fields; a sphere's distance function is one
  subtraction and one square root, the cheapest and most robust collider it
  has. For ``does the pipeline train'' the shape of the contact patch does
  not matter. Line 22--23 records that wheel positions and radii were kept
  from the first working version so the soil coupling did not move.
\lis{25}{26} A body whose geoms have zero total mass has a zero inertia
  tensor, and the solver divides by it. Every body here carries a geom with an
  explicit \code{mass}.
\lis{28}{29} The coordinate convention: $+x$ forward, $+z$ up, wheel bottoms
  at $z=0$ when the chassis origin is at $z=0.14$.
\end{explain}

\section{Imports and dimensions}

\codepart{\srcroot/luna_hifi_tasks/tricycle/tricycle.py}{31}{66}

\begin{explain}
\li{31} Lets the annotations on line 73 (\code{str | None},
  \code{tuple[float, float]}) be written in the modern form on any Python
  3.7+.
\lis{33}{35} \code{math} for the rail geometry, \code{os} and
  \code{tempfile} for \code{write_mjcf}.
\li{41} Sphere radius of all three wheels, 10~cm.
\li{42} Height of the chassis body origin above the ground when the wheels
  rest on $z=0$. \code{tricycle_env_cfg.py} imports this to compute the spawn
  height.
\li{43} Wheel centres relative to the chassis origin: $0.10 - 0.14 = -0.04$.
  Defined once so the three wheel bodies cannot drift apart.
\li{45} $x$ of the steering axis and the front wheel centre, 22~cm ahead of
  the chassis origin.
\li{46} $x$ of the rear axle, 15~cm behind. Wheelbase $= 0.37$~m.
\li{47} Half the rear track: rear wheel centres at $y=\pm0.16$.
\li{50} The apex of the triangle, directly over the steering axis.
\li{51} The rear corners of the triangle sit at $y=\pm0.13$, 3~cm inside the
  wheel centres, so the rails end inside the wheel spheres and the corners
  are visually capped by the wheels (Section~\ref{sec:new_body}).
\lis{52}{53} Rail cross-section: half-width 2~cm, half-height 1.5~cm, i.e.\
  4~cm wide and 3~cm tall. MuJoCo box sizes are always half-extents.
\lis{56}{58} Geom masses. $0.8 + 0.8 + 0.8 + 2.6 = 5.0$~kg, the same as the
  old single box, so the car's weight, wheel loads and drive torques are
  unchanged. Compiled with MuJoCo, the chassis centre of mass lands at
  $x=-0.054$~m (slightly behind the origin, towards the driven wheels) and
  the principal inertia is about $(0.062, 0.044, 0.019)$~kg\,m$^2$.
\lis{60}{62} Joint names. These strings appear in the MJCF, survive the USD
  conversion as prim names, and are what \code{tricycle_env.py} passes to
  \code{find_joints}. \code{JOINT_REAR} is a list because the two rear wheels
  are always addressed together.
\lis{64}{66} Colours as MJCF \code{rgba} strings. The converter copies them
  into \code{primvars:displayColor}, so they are what the Newton viewer
  shows.
\end{explain}

\section{Building the frame}

\codepart{\srcroot/luna_hifi_tasks/tricycle/tricycle.py}{68}{99}

\begin{explain}
\li{73} A helper that returns one \code{<geom>} element for a box whose local
  $x$ axis runs from \code{start} to \code{end} in the chassis $xy$-plane.
  Writing the rails this way means the numbers are derived from the corner
  coordinates rather than typed by hand.
\li{74} Docstring.
\li{75} The vector from start to end.
\li{76} Half the rail length; MuJoCo wants half-extents. For the rails,
  $\tfrac12\sqrt{0.37^2+0.13^2} = 0.1961$~m.
\li{77} The yaw angle of that vector. For the left rail (from $(-0.15,
  0.13)$ to $(0.22, 0)$) it is $\operatorname{atan2}(-0.13, 0.37) =
  -0.338$~rad $\approx -19.4^\circ$; the right rail gets $+0.338$.
\li{78} The rail's centre: the midpoint, $(0.035, \pm0.065)$.
\lis{79}{82} Emit the element. \code{pos} is the centre, \code{euler="0 0
  yaw"} rotates about $z$ (the compiler is told on line 107 that angles are
  radians), \code{size} is (half-length, half-width, half-height),
  \code{mass} fixes the mass (MuJoCo derives the density and the inertia),
  \code{rgba} the colour. The \code{name} lets the geom be found in the USD.
\li{85} The left rail: from the left rear corner to the apex.
\li{86} The right rail: mirror image.
\lis{89}{93} The crossbar, a box centred on the rear axle line
  ($x=-0.15$), half-length in $y$ of $0.13+0.02 = 0.15$ so that it covers
  the rails' ends, and the same cross-section as the rails.
\lis{96}{99} The deck: a $12\times13\times2$~cm plate centred at
  $x=-0.08$, i.e.\ filling the rear half of the triangle where it is wide
  enough, 5~mm lower than the rails' top surface so it reads as a floor pan.
  It carries 2.6 of the 5~kg.
\end{explain}

\begin{whybox}{why boxes and not a triangular mesh}
MJCF has no triangle primitive. A mesh with inline vertices would survive
MuJoCo's compiler, but boxes are the safest shape for every stage of the
pipeline: the USD converter writes them as \code{Cube} prims with a scale,
Newton's rigid collision pipeline handles boxes analytically, and the Isaac
Lab visualiser draws them without any mesh import. Three rails plus a deck
give an unmistakable triangle at zero risk.
\end{whybox}

\section{The MJCF string}

\codepart{\srcroot/luna_hifi_tasks/tricycle/tricycle.py}{101}{143}

\begin{explain}
\li{105} An f-string: every \code{{...}} is filled in by Python at import
  time, so the constants above and the rail strings are pasted into the XML.
\li{106} Root element; \code{model} is just a label.
\li{107} Compiler options. \code{angle="radian"} makes \code{euler} and joint
  \code{range} values radians (the steering range on line 123 is
  $\pm0.6$~rad). \code{autolimits="true"} tells MuJoCo that a joint with a
  \code{range} is limited without needing a separate \code{limited="true"}.
\lis{108}{111} Defaults inherited by every geom and joint. \code{condim="3"}
  gives contacts normal force plus two friction directions (sliding friction,
  no torsional/rolling terms). \code{friction="0.8 0.005 0.0001"} sets
  sliding, torsional and rolling friction coefficients; the converter copied
  these into the USD physics material (\code{physics:dynamicFriction = 0.8},
  \code{newton:torsionalFriction = 0.005}, \code{newton:rollingFriction =
  0.0001}). \code{joint damping="0.02"} adds a small viscous torque to every
  hinge so free wheels do not spin forever.
\li{113} Everything with a physical pose lives under \code{<worldbody>}.
\li{114} The chassis body, placed at $z=0.14$ in the world when the model is
  loaded on its own.
\li{115} A free joint: six degrees of freedom, so the chassis floats and can
  translate and rotate freely. Isaac Lab treats a body with a free joint as
  a \emph{floating-base} articulation, and its pose is the ``root state''
  that \code{tricycle_env.py} reads and writes.
\lis{116}{119} The four frame geoms, pasted in from lines 85--99.
\li{121} Comment.
\li{122} The fork body, a child of the chassis, at the front wheel centre.
  Its origin is the steering axis.
\li{123} The steering joint: a hinge about the fork's $z$ axis with range
  $\pm0.6$~rad ($\pm34.4^\circ$). The USD converter writes this as
  \code{physics:lowerLimit = -34.377} degrees. \code{tricycle_env_cfg.py}
  commands it with a position target of at most $\pm0.5$~rad, inside the
  limit.
\li{124} A thin vertical capsule (radius 1.2~cm, from $z=+0.02$ to
  $z=-0.02$ in the fork frame) so the fork body has mass (0.1~kg) and hence
  inertia. It sits entirely inside the front wheel sphere.
\li{125} The front wheel body, a child of the fork, at the fork origin.
\li{126} The wheel's spin axis: a hinge about $y$. No range, no actuator:
  the front wheel free-rolls on its default joint damping.
\li{127} The wheel sphere, 0.5~kg.
\lis{128}{129} Close the wheel and fork bodies.
\li{131} Comment.
\lis{132}{135} The left rear wheel: a body at $(-0.15, 0.16, -0.04)$ with a
  hinge about $y$ named \code{rear_left} and a 0.5~kg sphere. This joint is
  velocity-driven from the env.
\lis{136}{139} The right rear wheel, mirrored in $y$.
\lis{140}{143} Close chassis, worldbody, mujoco, and the f-string.
\end{explain}

\begin{isaacbox}{what Isaac Lab needs from this model}
Isaac Lab's \code{Articulation} class needs: exactly one root body (the
chassis, with its free joint), joints with stable names (for
\code{find_joints}), and every body with mass. It gets the joint
\emph{order} from the USD; \code{find_joints} returns indices in that order,
which is why the env never assumes a fixed index.
\end{isaacbox}

\section{Writing the file}

\codepart{\srcroot/luna_hifi_tasks/tricycle/tricycle.py}{146}{156}

\begin{explain}
\li{146} \code{write_mjcf} takes an optional directory and returns the path
  written.
\lis{147}{150} Docstring: idempotent, offline use only. The runbook's
  conversion step calls this function through
  \code{python -c "from luna_hifi_tasks.tricycle.tricycle import write_mjcf; print(write_mjcf())"}.
\li{151} Default location: the system temp directory plus
  \srcpath{luna_hifi_assets}. On Windows that is the folder the runbook's
  converter command reads from:
\begin{lstlisting}[style=shell]
C:\Users\<you>\AppData\Local\Temp\luna_hifi_assets\tricycle.xml
\end{lstlisting}
\li{152} Create the folder if needed.
\li{153} The file is always called \code{tricycle.xml}.
\lis{154}{155} Write the XML with the surrounding blank lines stripped and a
  final newline.
\li{156} Return the path so the caller can print it.
\end{explain}

\begin{whybox}{no more write-at-import}
The previous version of this file ended with
\code{TRICYCLE_MJCF_PATH = write_mjcf()}, so every \code{import} of the
package wrote a file into the temp directory, including on every training
start, and the constant was never used. The line has been removed; the
function is called explicitly when converting.
\end{whybox}
